\section{Introduction}

Koopman operator theory \citep{koopman_hamiltonian_1931, koopman_dynamical_1932, lan_linearization_2013,brunton_koopman_2016, brunton_modern_2022} offers a way to recast nonlinear dynamical systems as a linear dynamical system by working in a higher-dimensional (ideally infinite-dimensional) function or feature space rather than in the original state space.
Instead of evolving the state in the nonlinear dynamical system directly, Koopman theory studies the evolution of observables, i.e., functions (embeddings/features) of the original state, using a linear Koopman operator.
When one can find observables that approximately span a subspace that is invariant under the Koopman operator, the dynamics in that observable space are approximately linear, even if the original state dynamics are highly nonlinear (see Section~\ref{sec:background} for details).

However, global linearization is mathematically impossible for systems with multiple stable equilibria \citep{lan_linearization_2013, brunton_koopman_2016}.
% This implies that if a nonlinear system has more than one isolated fixed point, it cannot admit a finite-dimensional Koopman-invariant subspace that both (a) contains the state variables and (b) yields exact dynamics on all basins of attraction simultaneously.
In a single finite-dimensional global system with more than one isolated fixed point, we can nicely linearize the dynamics within each basin of attraction separately to admit classical linear-quadratic regulator (LQR) \citep{kalman_general_1960} or model predictive control (MPC) \citep{korda_linear_2018} style solutions, but the linearizations needed in different basins are incompatible with one another in a single finite-dimensional global system. The global nonlinear problem requires us to understand how trajectories cross basin boundaries and how to combine the locally linear subproblems.

% This observation creates a difficulty for control.
In many control problems, the task is precisely to move between different basins of attraction: for example, to drive the system from one stable operating point to another, or to traverse a state space that is partitioned into different dynamical regimes. 
So, the crux of the problem is this: Given a nonlinear control problem which can be partitioned into basins of attraction, where we can solve each basin's local dynamics in closed form, \textbf{how do we model movement from one basin of attraction to another?}
How do we traverse from one side of the field to another by passing through the basins of attraction?
Can we represent this nonlinear dynamical system like a graph, where we can construct an adjacency matrix and traverse a graph where we solve each node (basin of attraction) locally?
This is essentially about how to \textbf{combine} the solutions to smaller linear problems to solve a hard global nonlinear control problem.

Our proposed approach marries two simple ideas: periodic re-encoding and sparse coding.
At a high level, periodic re-encoding sustains the quality of the latent embedding over time by course-correcting the latent state as trajectories move across regions of state space. Then, representing the latent as a sparse code induces a union-of-subspaces structure that should naturally separate and select the correct local linear Koopman dynamics.
The support pattern of the sparse code acts like a mode label indicating which local linear model is currently active, and transitions between supports define a graph over modes.
Planning then becomes a graph traversal problem over these sparse-code regions, while each node carries its own local linear dynamical model that can be leveraged with LQR- or Koopman-MPC-style tools.

\section{Background}
\label{sec:background}

\subsection{Koopman operator theory}
\label{sec:koopman-background}

This subsection largely follows the review of \citet{brunton_modern_2022} with some changes in notation. \note{It's meant to be a bit more accessible; I haven't taken measure theory. Sorry.}

\paragraph{Koopman operator and observables.}
Consider a discrete-time nonlinear dynamical system
\begin{equation}
x_{t+1} = f(x_t), \qquad x_t \in \mathcal{X} \subset \mathbb{R}^{d_x},
\end{equation}
where $x_t$ is the state and $f : \mathcal{X} \to \mathcal{X}$ describes the dynamics.
Koopman operator theory studies the evolution of \emph{observables}, which are measurement functions $h : \mathcal{X} \to \mathbb{R}$.
The Koopman operator $\mathcal{K}$ is defined by composition:
\begin{equation}
(\mathcal{K}h)(x) = h(f(x)).
\end{equation}
Although $f$ may be nonlinear, $\mathcal{K}$ is always linear on the (typically infinite-dimensional) function space of observables. Then, instead of evolving the state $x_t$ directly, we can evolve a vector of observables
\begin{equation}
\phi(x) = \big[h_1(x), \dots, h_{d_z}(x)\big]^\top.
\end{equation}
If the span of $\{h_1,\dots,h_{d_z}\}$ is invariant under $\mathcal{K}$, 
then there exists a matrix
$K \in \mathbb{R}^{d_z \times d_z}$ such that
\begin{equation}
\phi(x_{t+1}) = \phi(f(x_t)) = \mathcal{K}\phi(x_t) = K\, \phi(x_t),
\end{equation}
so that the nonlinear dynamics of $x_t$ are exactly represented by a linear dynamical system
in the observable coordinates $z_t = \phi(x_t)$.
To learn linear latent-state models of nonlinear dynamics, the goal is to find observables that approximately span such a Koopman-invariant subspace.

\paragraph{Approximating the Koopman operator: DMD/eDMD.}
In practice, we do not have access to the infinite-dimensional Koopman operator, and we must work with a finite family of observables.
Given a feature map $\phi : \mathcal{X} \to \mathbb{R}^{d_z}$, one seeks a finite-dimensional approximation
\begin{equation}
\phi(x_{t+1}) \approx K\, \phi(x_t),
\end{equation}
where $K \in \mathbb{R}^{d_z \times d_z}$ is a learned linear model.
Data-driven approximations such as dynamic mode decomposition (DMD) \citep{schmid_dmd_2010, rowley_2009_spectral, tu_dynamic_2014} and extended DMD (eDMD) \citep{williams_datadriven_2015, williams_kernel-based_2016} estimate $K$ from pairs $\{(x_t, x_{t+1})\}$ by solving
\begin{equation}
K = \arg\min_{\hat K} \sum_{t} \big\| \phi(x_{t+1}) - \hat K\, \phi(x_t) \big\|_2^2,
\end{equation}
with $\phi$ chosen from a (possibly large) dictionary of basis functions.
If the chosen observables are approximately invariant under $\mathcal{K}$, then the learned $K$ captures a
nearly linear evolution of the feature vector $z_t = \phi(x_t)$ while still allowing rich nonlinear behaviour
in the original state $x_t$.

DMD is a special case: we take $\phi(x)=x$ itself, so DMD directly fits a linear map
\begin{equation}
K^{\mathrm{DMD}} = \arg\min_{\hat K} \sum_{t} \big\| x_{t+1} - \hat K x_t \big\|_2^2
\end{equation}
to the state.
This can be viewed as learning a representation in which the transition model is as close as possible to linear while still capturing important nonlinearity in the observed trajectories.


% \paragraph{Dynamical systems definitions.}
% Recall the key notions from nonlinear dynamics that will be important for our multi-basin problem.



% A \emph{non-wandering set} $S$ is a set of points with the property that for every neighborhood $U$ of any $x \in S$ and every $T > 0$, there exists a time $t \ge T$ and an initial condition $x_0 \in U$ such that the trajectory returns to $U$ at time $t$.
% Informally, trajectories starting in $S$ keep returning arbitrarily close to $S$.

% An \emph{attractor} $A$ is a non-wandering set that has the property of asymptotic stability: trajectories that start near the set approach it asymptotically as $t \to \infty$. The attractor has an open \emph{basin of attraction}
% \begin{equation}
% \mathcal{B}(A) = \big\{ x : \operatorname{dist}(\phi_t(x), A) \to 0 \text{ as } t \to \infty \big\},
% \end{equation}
% where $\phi_t(x)$ denotes the trajectory at time $t$ starting from $x$.
% The basin of attraction is just the set of initial conditions whose trajectories approach the attractor $A$ in the long-time limit.

\paragraph{Problems with multiple basins of attraction.}

While Koopman representations are powerful, they face theoretical limitations in systems with complex global topology. A key concept here is the \emph{basin of attraction}. First, an \emph{attractor} $A$ is a set of points that has the property of asymptotic stability: trajectories that start near the set approach it asymptotically as $t \to \infty$ 
(e.g. a \emph{fixed point} $x^\ast$ that satisfies $f(x^\ast) = x^\ast$). Then, the basin of attraction $\mathcal{B}(A)$ is the set of initial conditions whose trajectories approach the attractor $A$ in the long-time limit:
\begin{equation} 
\mathcal{B}(A) = \{ x \in \mathcal{X} : \lim_{t \to \infty} \text{dist}(f^t(x), A) = 0 \}.
\end{equation}

For many nonlinear systems, especially those with multiple basins of attraction, no finite-dimensional Koopman-invariant subspace exists that both (1) contains the original state coordinates $x$, and (2) is globally invariant under $\mathcal{K}$ in a way that yields exact dynamics on all basins simultaneously.

\citet{lan_linearization_2013}, also summarized in \citet{brunton_koopman_2016}, show that such finite-dimensional invariant subspaces containing the state can exist only for systems with a single isolated fixed point whose basin covers the state space.
% If a nonlinear system has more than one isolated fixed point, it cannot admit a finite-dimensional Koopman-invariant subspace that both spans the state and yields exact global dynamics.
% However, for each \emph{single} basin of attraction, we can linearize the dynamics within that basin nicely.
% Near a hyperbolic fixed point, classical linearization theory guarantees the existence of local coordinates in which the dynamics are well-approximated by a linear system, and Koopman eigenfunctions provide one way to construct such local coordinates.
% The only struggle is: how do we glue these separate basins, where we can linearize the dynamics separately, into one finite-dimensional global system that covers all basins and still keeps the dynamics linear and the map from latent $z$ to state $x$?
Any kind of finite-dimensional Koopman-learned latent for a multi-basin system must either:
\begin{itemize}
    \item be effectively \emph{piecewise linear}, where the latent includes some implicit ``which basin am I in?'' variable, or
    \item be \emph{approximate}, in the sense that it smooths over the boundaries between regions with different dynamics and is not exact anywhere.
\end{itemize}
This observation motivates our use of sparse coding to make explicit which local model is active, and periodic re-encoding to switch between locally valid linearizations as trajectories move across basin boundaries.



\paragraph{Koopman predictors for controlled systems (optional).}
For controlled systems
\begin{equation}
x_{t+1} = f(x_t, u_t),
\end{equation}
Koopman-based approaches often add controls to the higher-dimensional latent where the dynamics are linear:
\begin{equation}
z_{t+1} = A z_t + B u_t, \qquad x_t \approx C z_t,
\end{equation}
where $z_t = \psi(x_t)$ is a vector of learned observables, and $A,B,C$ are learned from data \citep{korda_linear_2018}.
This predictor has linear dynamics in the lifted space, but the feedback law $u_t = \kappa_{\mathrm{lift}}(z_t)$ induces a nonlinear controller $\kappa(x_t) = \kappa_{\mathrm{lift}}(\psi(x_t))$ in the original state.
Such linear predictors can be plugged into standard tools such as LQR and MPC while the nonlinearity is hidden in the fixed lifting map $x \mapsto z$ and decoder.

There are many benefits to having linear dynamics in the latent space.
Firstly, linear dynamics allow for LQR solvers to be applied to derive closed-form solutions to optimal control problems \citep{kalman_general_1960}.
Linear dynamics also significantly simplify system identification and interpretability, and can be efficiently unrolled with parallelism (e.g., in modern state-space models \citep{gu2022efficiently, gu_2023_mamba, smith2023simplified, dao_2024_mamba2}), which is attractive for long-horizon prediction and planning.



\subsection{Koopman autoencoders and periodic re-encoding}
\label{sec:background-reencoding}

Recent work \citep{lusch_deep_2018, azencot_forecasting_2020, shi_deep_2022, frion_leveraging_2023} has combined Koopman operator theory with deep autoencoders to learn nonlinear observables and linear latent dynamics directly from data, similar to DMD and eDMD.
In a typical Koopman autoencoder, an encoder $\phi : \R^{d_x} \to \R^{d_z}$ and decoder $\psi : \R^{d_z} \to \R^{d_x}$ are trained together with a linear latent Koopman transition $K$ so that
\begin{equation}
z_t \approx \phi(x_t), \qquad z_{t+1} \approx K z_t, \qquad x_t \approx \psi(z_t).
\end{equation}
We might do this by minimizing the following objective:
\begin{equation}
\min_{K, \phi, \psi} \sum_\mathcal{D} ||x_t - \psi(\phi(x_t)) ||_2 + \lambda \cdot || \phi(x_{t+1}) - K \phi(x_t)||_2
\end{equation} 
where $\lambda > 0$ is a scalar.
Once trained, one can simulate forward in time by iterating the linear latent dynamics and decoding:
\begin{equation}
z_{t+k} \approx K^k z_t, \qquad \hat x_{t+k} = \psi(z_{t+k}).
\end{equation}

\citet{fathi_course_2023} observed that unrolling the linear dynamics in latent space in this way leads to long-term drift in trajectories when mapped back to the state space.
Because the encoder and decoder are not exact inverses and the Koopman-invariance of the learned features is only approximate, latent trajectories $z_t$ can slowly drift away from the region where the local linear approximation is valid, resulting in trajectories that cross or develop spurious attractors when decoded to $x$.

To address this, they introduced \emph{periodic re-encoding}.
Rather than letting the latent state evolve indefinitely under $K$, one intermittently decodes and re-encodes:
\begin{equation}
z_{t+k} = \phi\!\left(\psi\!\left(K^k z_t\right)\right),
\end{equation}
for some re-encoding period $k$.
More concretely, the trajectory is evolved linearly in the latent space for several steps, decoded back to the original state, and then this decoded state is re-encoded to obtain a fresh latent representation.
This contrasts with a one-shot encoding at the initial time, where the same linear approximation is used regardless of how far the trajectory has moved in state space.

In the context of the result by \citet{lan_linearization_2013}, periodic re-encoding can be understood as a mechanism for switching between locally-valid linearizations as the dynamics traverse different regions or cross basin boundaries.
Within a single basin away from the boundaries, the encoder–decoder pair should behave approximately linearly along trajectories, but near basin boundaries, the geometry of invariant sets and eigenfunctions changes rapidly: observables that were approximately Koopman-invariant in one region may no longer span an invariant subspace in the new region.
By re-encoding at these times (or at a fixed period that is short compared to the time spent near boundaries), the model effectively refreshes the latent embedding so that the latent linear dynamics are again centered on the current part of the attractor or basin. This should help the system transition cleanly between different basins of attraction.

\subsection{Sparse coding}
\label{sec:background-sparse}

Sparse coding seeks a representation $z$ such that $x \approx Dz$ for a learned dictionary $D : \R^{d_z} \to \R^{d_x}$, via the optimization
\begin{equation}
\min_z \;\; \frac{1}{2} \|x - D z\|_2^2 + \lambda \|z\|_1,
\end{equation}
where $\lambda > 0$ controls the sparsity of $z$ \citep{olshausen_emergence_1996, chen_atomic_1998, donoho_optimally_2003}.
The $\ell_1$ penalty induces a mapping $x \mapsto z(x)$ that is piecewise linear and piecewise constant in its support: on each region of input space, the same small set of dictionary atoms is active, and within that region, $z(x)$ depends linearly on $x$.
Geometrically, this models the latent space as a union of lower-dimensional linear subspaces, with each subspace indexed by the support (active set) of the sparse code.

At a high level, this union-of-subspaces structure is exactly the kind of structure we would like to impose on a multi-basin nonlinear control problem.
Each support pattern corresponds to a locally linear regime; moving between regimes corresponds to changing which atoms are active.

In our setting, we want to treat the Koopman features as a sparse combination of dictionary atoms.
The dictionary $D$ is global, and the support pattern of the sparse code $z$ acts like a mode label indicating which basin or local region we are currently in.
Inside a region (basin of attraction) where the dynamics look locally linear, that region uses the same dictionary atoms and therefore has its own local Koopman-like dynamics.
Crossing a basin boundary is equivalent to changing the sparse code support to another region with a different local linear model. I felt that this was also explained nicely in Appendix D of \citet{fathi_course_2023}.

This gives a natural way to build an adjacency matrix for a graph that encapsulates the nonlinear problem: nodes correspond to local regions of constant support, and edges correspond to observed transitions between supports along trajectories.
On each node, one can learn a local Koopman linear operator and solve the associated LQR-style control problem in closed form.
Planning trajectories from one side of the field to the other then becomes a graph traversal problem over these sparse-code regions.

\paragraph{LISTA as a fast sparse encoder.}
Classical sparse coding solves the above optimization with an iterative algorithm (e.g., ISTA/FISTA, coordinate descent), which is often too slow to use inside a learned dynamical model.
LISTA (Learned ISTA) \citep{gregor_learning_2010} replaces the iterative solver with a trainable, finite-depth network that approximates the sparse solution.

In LISTA, one learns a non-linear encoder $z = f_\theta(x)$ whose architecture mimics a truncated iterative shrinkage/thresholding algorithm.
The key nonlinearity is a component-wise shrinkage function with learned thresholds; by construction, the mapping $x \mapsto z(x)$ is explicitly piecewise affine, with regions defined by which coordinates survive thresholding.
This builds a strong inductive bias for \emph{mode selection} into the encoder: for each $x$, only a few coordinates of $z$ are active, picking a small subset of dictionary atoms that explain the input.
Compared to a generic nonlinear encoder where all latent coordinates can participate everywhere, LISTA-style encoders encourage a sparse, structured partition of state space into locally linear regimes, which is exactly what we will leverage in our Koopman control architecture.


\begin{algorithm}[ht]
\caption{LISTA Encoder (Forward Pass)}
\label{alg:lista}
\begin{algorithmic}[1]
\REQUIRE $x \in \mathbb{R}^{d_x}$; depth $T \in \mathbb{N}$; encoder $\phi$ with learnable parameters $W_e \in \mathbb{R}^{d_z \times d_x}$, $S \in \mathbb{R}^{d_z \times d_z}$, $\theta \in \mathbb{R}^{d_z}_{\ge 0}$
\ENSURE Output $z \in \mathbb{R}^{d_z}$ ($z$ is latent code $\phi(x)$)
\STATE $z^{(0)} \gets h_\theta (W_e x)$
\FOR{$t = 1$ \TO $T$}
    \STATE $c^{(t)} \gets W_e x + Sz^{(t-1)}$ 
    \STATE $z^{(t)} \gets \operatorname{h}_{\theta} \big(c^{(t)}\big)$
\ENDFOR
\RETURN $z^{(T)} = \phi(x)$
\end{algorithmic}
\end{algorithm}

We represent the decoder as a linear dictionary $D \in \mathbb{R}^{d_x \times d_z}$, so that $\hat{x} = D z$. The encoder $\phi(x)$ is implemented as a LISTA network (Alg. \ref{alg:lista}) with parameters $W_e \in \mathbb{R}^{d_z \times d_x}, S \in \mathbb{R}^{d_z \times d_z}$, and thresholds $\theta \in \mathbb{R}^{d_z}_{\ge 0}$; here $W_e$ plays the role of a learned operator analogous to $D^\top$, while $D$ is the dictionary used in reconstruction.

\section{Method}
\label{sec:method}

We propose a method to learn a global Koopman operator for nonlinear dynamical systems characterized by multiple basins of attraction. Our key insight is to enforce a \textit{sparse} latent representation, structured as a union of subspaces, where each active support approximately corresponds to a specific basin of attraction or dynamical regime. By learning a sparse encoder and a global linear transition operator, we effectively partition the state space into regions where local linear dynamics hold, while maintaining a consistent global latent space.

\subsection{Problem Formulation}
Consider a discrete-time dynamical system $x_{t+1} = F(x_t)$ with state $x_t \in \mathbb{R}^{d_x}$. We seek an embedding function $\phi: \mathbb{R}^{d_x} \to \mathbb{R}^{d_z}$ and a linear operator $K \in \mathbb{R}^{d_z \times d_z}$ such that the dynamics in the latent space are linear: $\phi(x_{t+1}) \approx K \phi(x_t)$. To map back to the state space, we learn a decoder $\psi: \mathbb{R}^{d_z} \to \mathbb{R}^{d_x}$ such that $x_t \approx \psi(\phi(x_t))$.

\subsection{Sparse Koopman Autoencoder}
Our architecture is a deep autoencoder regularized to learn Koopman invariant subspaces that are sparse. 

\paragraph{Generic Sparse Architecture.}
As a baseline, we employ a standard Multi-Layer Perceptron (MLP) for both the encoder $\phi_\theta$ and decoder $\psi_\gamma$. The latent state $z_t = \phi_\theta(x_t)$ is regularized with an $\ell_1$ penalty to induce sparsity. The decoder reconstructs the state as $\hat{x}_t = \psi_\gamma(z_t)$.

\paragraph{LISTA Encoder.}
To enforce structured sparsity and interpretability, we replace the generic MLP encoder with a Learned Iterative Soft-Thresholding Algorithm (LISTA) network \citep{gregor_learning_2010}. LISTA is an unrolled optimization algorithm that approximates the solution to the sparse coding problem:
\begin{equation}
    z^* = \argmin_z \frac{1}{2} \|x - D z\|_2^2 + \lambda \|z\|_1,
\end{equation}
where $D \in \mathbb{R}^{d_x \times d_z}$ is the decoder dictionary.
Our encoder $\phi(x)$ consists of $T$ iterations of the update rule:
\begin{equation}
    z^{(t+1)} = \mathcal{S}_{\alpha} (W_e x + S z^{(t)}),
\end{equation}
where $\mathcal{S}_\alpha(\cdot)$ is the element-wise soft-thresholding operator with threshold $\alpha$, and $W_e$ and $S$ are learnable weight matrices. In our implementation, we optionally preprocess $x$ with a shallow MLP before feeding it into the LISTA layers (referred to as \textit{Nonlinear LISTA}) to allow for a nonlinear transform into the sparse basis. The decoder is strictly linear, $\psi(z) = D z$, where the columns of $D$ serve as the dictionary atoms.

\paragraph{HyperLISTA Module.}
We also explore a parameter-efficient variant called HyperLISTA \citep{chen_hyperparameter_2021}, which couples the encoder and decoder weights. Instead of learning $W_e$ and $S$ independently, HyperLISTA derives them analytically from the current decoder dictionary $D$:
\begin{equation}
    W_e = \frac{1}{L} D^\top, \quad S = I - \frac{1}{L} D^\top D,
\end{equation}
where $L$ is an estimate of the Lipschitz constant (spectral norm) of $D^\top D$. This formulation reduces the trainable parameters of the encoder to a small set of scalars regulating the thresholding dynamics (e.g., momentum and support selection), ensuring that the encoder update is always consistent with the sparse coding objective for the current dictionary.

\subsection{Learning Objective}
We train the parameters of the encoder, decoder, and the Koopman matrix $K$ jointly by minimizing a composite loss function comprising four terms:
\begin{equation}
    \mathcal{L} = \lambda_{\text{recon}} \mathcal{L}_{\text{recon}} + \lambda_{\text{pred}} \mathcal{L}_{\text{pred}} + \lambda_{\text{align}} \mathcal{L}_{\text{align}} + \lambda_{\text{sparse}} \mathcal{L}_{\text{sparse}}.
\end{equation}
\begin{itemize}
    \item \textbf{Reconstruction Loss}: Ensures the autoencoder faithfully represents the state, $\mathcal{L}_{\text{recon}} = \|x_t - \psi(\phi(x_t))\|_2^2$.
    \item \textbf{Prediction Loss}: Enforces accurate forward prediction in the state space via the latent dynamics, $\mathcal{L}_{\text{pred}} = \|x_{t+1} - \psi(K \phi(x_t))\|_2^2$.
    \item \textbf{Alignment (Residual) Loss}: Enforces linearity directly in the latent space, $\mathcal{L}_{\text{align}} = \|\phi(x_{t+1}) - K \phi(x_t)\|_2^2$.
    \item \textbf{Sparsity Loss}: Induces sparsity in the code $z$, $\mathcal{L}_{\text{sparse}} = \|\phi(x_t)\|_1$.
\end{itemize}
Spectral normalization and homogeneous coordinates are optionally applied to stabilize training.

\subsection{Inference with Periodic Re-encoding}
Long-term forecasting with purely linear latent dynamics ($z_{t+k} = K^k z_t$) often suffers from drift, where the latent state typically departs the valid manifold of the encoder. To mitigate this, we employ \textit{periodic re-encoding} \citep{fathi_course_2023}. Every $M$ steps, we project the predicted latent state back to the physical state space and re-encode it:
\begin{equation}
    z_{t+M} \leftarrow \phi(\psi(z_{t+M})).
\end{equation}
This mechanism acts as a "correction" step, re-centering the latent state within the appropriate basin of attraction and allowing the sparse encoder to switch active sub-dictionaries as the trajectory crosses dynamical boundaries.

\section{Experiments}
\label{sec:experiments}

We evaluate our method on a suite of dynamical systems exhibiting multiple basins of attraction and chaotic behavior. We compare the standard MLP-based Koopman Autoencoder (\textbf{GenericKM}) against the LISTA-based architectures (\textbf{LISTAKM} and \textbf{HyperLISTAKM}).

\subsection{Environments}

\paragraph{Duffing Oscillator (2D).}
A classic nonlinear oscillator with a double-well potential, defined by $\ddot{x} + \delta \dot{x} - \alpha x + \beta x^3 = 0$. The system features two stable spirals and a saddle point at the origin. This serves as a canonical testbed for separating distinct basins of attraction.

\paragraph{Lyapunov Multi-Attractor System.}
To test the capacity limit of the sparse code, we construct a generic gradient field system with $N=13$ exponentially stable fixed points arranged in a grid. The dynamics are governed by $\dot{x} = -\nabla V(x)$, where $V(x)$ is constructed from a superposition of Gaussian wells. This environment explicitly tests the model's ability to isolate local linearizations for many distinct attractors using a single global model.

\paragraph{Multi-Basin Chaotic Systems (Dysts).}
We extend our evaluation to a curated subset of the \texttt{dysts} benchmark database \citep{gilpin_dysts_2021}, focusing on high-dimensional chaotic systems that exhibit multi-scroll attractors or multiple coexisting attractors (e.g., the Lorenz, Chen, and Chua, and related multi-wing systems). These systems represent a significantly harder challenge due to their chaotic nature and complex attractor topology.

\subsection{Experimental Setup}
Models are trained on trajectory pairs $(x_t, x_{t+1})$ sampled uniformly from the state space of interest. We evaluate forecasting performance using Mean Squared Error (MSE) over long-horizon rollouts. To rigorously test stability, we compare three inference modes: (1) \textit{Open-loop}, where the latent state evolves linearly forever; (2) \textit{Every-step re-encoding}, where the state is decoded and re-encoded at every step; and (3) \textit{Periodic re-encoding}, with periods $M \in \{10, 25, 50, 100\}$.

\subsection{Results}
Our experiments aim to demonstrate that the structured sparsity induced by LISTA allows for more stable long-term predictions compared to dense encodings, particularly when coupled with periodic re-encoding to manage transitions between local linear regimes. We report quantitative MSE scores at horizons $T=100$ to $T=1000$ and provide qualitative phase portraits showing the model's ability to capture the global topology of the attractors.
